\chapter{条件核心：雅化集团（002497）}

\begin{dashboardbox}[结论卡]
\textbf{现价：}16.79元 \quad \textbf{总市值：}193.52亿元 \quad \textbf{正式区间：}约25--26元 \quad \textbf{牛市锚：}42元 \\
\textbf{基准/最终空间：}约+51\% \quad \textbf{牛市空间：}+150.1\% \quad \textbf{判断性翻倍概率：}20\% \\
\textbf{核心判断：}证据最完整的锂盐高弹性标的，民爆提供利润底仓；但FY2025和2026Q1经营现金流均为负，H2必须证明盈利不是库存与套保的短期组合。
\end{dashboardbox}

\section{为什么它是唯一条件核心}

雅化集团的优势不只是锂价弹性，而是双主业可以被分开验证。2025年锂业务收入48.24亿元，占总收入56.46\%，毛利率12.22\%；民爆业务收入33.16亿元，占38.82\%，整体毛利率37.46\%。锂业务决定上行斜率，民爆业务提供相对稳定的利润底盘，两者不能套用同一倍数。

公司年报披露当前锂盐综合产能近13万吨，2025年锂产品销量6.95万吨；长期合作客户包括Tesla、LGES、SK ON、宁德时代等，重大合同处于正常履行状态。头部客户收入占锂业务90\%以上。这些信息比仅有行业需求或客户认证的标的更接近公司级收入证据。

民爆业务同样有真实经营壁垒：牌照、区域渠道、项目服务和安全运营构成进入门槛。2025年工业炸药许可产能26.25万吨、利用率83.74\%；2026年4月许可产能增至33万吨。新增许可不等于新增利润，仍需观察利用率与回款。

\begin{exhibitbox}[Exhibit 4：双主业经营桥]
\small
\begin{tabularx}{\exhibitboxwidth}{L{2.3cm} L{3.2cm} L{3.2cm} X}
\toprule
分部 & 已披露经营证据 & 2026年关键预测/推断 & 估值处理 \\
\midrule
锂盐与资源 & 近13万吨产能；2025销量6.95万吨；头部客户与重大合同正常履行 & 券商预计全年约12万吨出货、自供比例约25\%；Q2吨利2.6--3.4万元 & 周期PE；剔除库存与套保一次影响 \\
民爆 & 2025收入33.16亿元、毛利率37.46\%；炸药利用率83.74\% & 券商预计利润约5.5--6亿元；新增许可产能需验证利用率 & 稳定分部信用，不给锂盐高弹性倍数 \\
运输与其他 & 2025收入4.03亿元 & 不作为核心增长来源 & 不单独给溢价 \\
\bottomrule
\end{tabularx}
\sourcenote{公司2025年报、2026H1业绩预告；东吴证券2026年7月7日原始报告。}
\end{exhibitbox}

\section{盈利从预告到全年还差什么}

公司2025年收入85.43亿元、归母净利润6.32亿元；2026Q1收入28.30亿元、归母3.39亿元、扣非3.57亿元。H1预告归母净利润11--13亿元、扣非11.25--13.15亿元。与券商2026全年30.32亿元归母净利润相比，H2仍需贡献17.32--19.32亿元。

这个H2要求不是简单的季节性外推，而是量、价、自供矿、库存与套保共同作用。原始报告预计2026年锂盐出货约12万吨、权益资源自供比例约25\%，并认为套保拖累在下半年减弱。公司公告只确认H1销量与销售均价同步增长，没有披露吨位、实现ASP与单位成本；因此吨利和出货必须继续标记为券商预测。

最关键的反证来自现金流。FY2025经营现金流为-5.70亿元，2026Q1为-4.31亿元。券商还预计2026年存货可能从16.96亿元增至51.13亿元。若利润增长主要来自库存价差，而现金继续被营运资金占用，则低PE并不代表可分配价值。

\begin{exhibitbox}[Exhibit 5：盈利—现金验证表]
\small
\begin{tabularx}{\exhibitboxwidth}{L{2.5cm} R{2cm} R{2cm} R{2cm} X}
\toprule
指标 & FY2025 & 2026Q1 & 2026H1预告 & 研究含义 \\
\midrule
营业收入 & 85.43亿元 & 28.30亿元 & 未披露 & 不从利润倒推收入 \\
归母净利润 & 6.32亿元 & 3.39亿元 & 11--13亿元 & H2仍需17.3--19.3亿元达到券商全年预测 \\
扣非净利润 & 6.94亿元 & 3.57亿元 & 11.25--13.15亿元 & 扣非改善较纯投资收益更可信 \\
经营现金流 & -5.70亿元 & -4.31亿元 & 未披露 & 核心阻断项，要求H2明显转正 \\
\bottomrule
\end{tabularx}
\sourcenote{公司2025年报、2026Q1报告、2026H1预告。H1预告未经审计。}
\end{exhibitbox}

\section{估值：为什么42元只属于牛市情景}

原始券商模型用2026E EPS 2.63元乘16倍PE，得到42.08元并给出42元目标。现价只对应6.38倍。若到年末翻倍至33.58元，需要12.77倍2026E PE，低于券商16倍但约为当前倍数的两倍。

本报告不让三档情景共用同一个EPS。悲观/基准/牛市分别假设锂盐出货9/12/13万吨、归母净利润13.1/30.35/38.5亿元、EPS 1.14/2.63/3.34元，再给予8/10/12.6倍PE，对应9.12/26.30/42.08元。概率分别为30\%/50\%/20\%，经营驱动概率值24.30元；再与现价和期限未披露的42元外部锚按80\%/10\%/10\%组合，审计值为25.32元，读者口径约25--26元。该方法明确表达：盈利兑现可以带来较大修复，但翻倍还需要更强盈利与倍数同时出现。

\begin{riskbox}[降级条件]
若Q3经营现金流继续明显为负、利润主要来自库存收益或套保、锂盐出货与自供比例未兑现，或民爆回款与利润恶化，雅化集团应从条件核心降为周期观察。价格继续下跌不是自动升级理由。
\end{riskbox}

\section{年末前的必要证据}

三季报需要看到锂盐出货接近3万吨、剔除库存与套保后吨利至少约2万元、经营现金流转正或显著改善；全年需要接近12万吨出货、归母利润接近30亿元、现金流转正、民爆利润约5.5--6亿元。吨位与吨利来自券商估计，现金阈值属于AStock验证要求。

只有当这些证据共同出现，市场把2026E EPS从6倍重估到13倍以上才具备基本面基础。若只出现锂价上涨而现金与自供比例不改善，42元仍只是周期尾部。

\clearpage
